\documentclass[11pt]{article}

\usepackage[margin=1in]{geometry}
\usepackage{amsmath,amssymb,amsthm,mathtools}
\usepackage{booktabs,longtable,array}
\usepackage[T1]{fontenc}
\usepackage{lmodern}
\usepackage{microtype}
\usepackage{xcolor}
\usepackage{hyperref}
\usepackage{xurl}

\hypersetup{
  colorlinks=true,
  linkcolor=blue!55!black,
  citecolor=blue!55!black,
  urlcolor=blue!55!black,
  pdftitle={A length-ten base for the rank-four counterexample: structure and deformation theory}
}

\newtheorem{theorem}{Theorem}[section]
\newtheorem{proposition}[theorem]{Proposition}
\newtheorem{lemma}[theorem]{Lemma}
\newtheorem{corollary}[theorem]{Corollary}
\theoremstyle{remark}
\newtheorem{remark}[theorem]{Remark}

\newcommand{\F}{\mathbf F}
\newcommand{\Z}{\mathbf Z}
\newcommand{\Spec}{\operatorname{Spec}}
\newcommand{\Soc}{\operatorname{Soc}}
\newcommand{\gr}{\operatorname{gr}}
\newcommand{\length}{\operatorname{length}}
\newcommand{\Ann}{\operatorname{Ann}}
\newcommand{\GL}{\operatorname{GL}}
\newcommand{\Ad}{\operatorname{Ad}}
\newcommand{\id}{\operatorname{id}}
\newcommand{\eps}{\varepsilon}
\newcommand{\m}{\mathfrak m}
\newcommand{\n}{\mathfrak n}

\title{A Length-Ten Base for the Rank-Four Counterexample:\\
Its Linear, Subgroup, and Deformation-Theoretic Structure}
\author{Exact compression and independent audit of the universal cubic construction}
\date{11 July 2026}

\begin{document}
\maketitle

\begin{abstract}
The previously audited rank-four group scheme not killed by four was defined
over a universal Artin ring of length \(632\).  This note compresses it to an
explicit characteristic-four local Gorenstein ring \(R\) of length \(10\)
and cardinality \(2^{10}=1024\).  We give a four-generator presentation of
\(R\), the complete multiplication and coproduct tables of the rank-four
Hopf algebra, and an explicit antipode.  Direct calculation gives
\[
 [4]^\#(e_1)=(2z)e_3\ne0,
 \qquad [4]^\#(e_2)=[4]^\#(e_3)=0,
 \qquad [8]^\#=\iota\eps.
\]
The construction is governed by the divided-power cubic
\[
 X^{[2]}Z+XYD+YD^{[2]}+Z^{[3]}.
\]
Conceptually, the fourth-power obstruction is the product of a
mixed-characteristic carry with the first noncocommutative symbol.  If
\(P=\iota\eps\), where \(\iota:R\to A\) is the unit, \(\phi=[2]^\#\), and
\(\psi=\phi-P\), then already on the
universal cubic base one has the exact identity
\[
 \psi^2=2\psi.
\]
Thus the cubic obstruction is intrinsically ``Bockstein-type carry times
noncommutativity.''  A bounded exact Boolean MinRank computation returns
UNSAT in rank at most three; conditional on that solver verdict, length
\(10\) is minimal among all Artin quotients of the audited universal cubic
base on which either fourth-power defect survives.  This is not asserted to
be a global minimum among every conceivable unrelated construction.  We
also describe the resulting group in several complementary ways.  Its
coordinate algebra is a two-equation relative complete intersection, its
dual is a four-dimensional PBW-type measure algebra, and reduced right
translation gives a closed immersion \(G\hookrightarrow\GL_{3,R}\), of
minimal possible dimension.  The group has exactly sixteen rank-two finite
flat subgroups over \(R\), all abstractly the same Oort--Tate group and none
normal.  Finally, a length-\(6\to8\to10\) deformation tower separates the two
phenomena: the first square-zero lift destroys a rank-two normal filtration,
while the last square-zero lift creates the nonzero fourth-power class.
\end{abstract}

\tableofcontents

\section{Statement of the result and scope}

The main construction is completely explicit.

\begin{theorem}\label{thm:main}
There is a finite local ring \(R\) and a finite locally free group scheme
\(G/\Spec R\) with the following properties.
\begin{enumerate}
\item \(R\) has characteristic \(4\), residue field \(\F_2\), cardinality
  \(2^{10}\), Hilbert function \((1,4,4,1)\), and one-dimensional socle.
\item \(G\) has rank \(4\), its unique geometric fiber is
  \(\alpha_2\times\alpha_2\), and \(G\) is noncommutative.
\item The fourth-power word is nontrivial, while the eighth-power word is
  trivial.  Thus \(G\) is not killed by its order:
  \[
    [4]\ne e,\qquad [8]=e.
  \]
\item \(G\) embeds as a closed subgroup of \(\GL_{3,R}\), but not of
  \(\GL_{2,R}\).  It has rank-two finite flat closed subgroups but no normal
  one, and hence no nontrivial finite-flat composition filtration.
\end{enumerate}
\end{theorem}

The earlier universal base had Hilbert function \((1,15,107,509)\), length
\(632\), and cardinality \(2^{632}\).  The new base therefore changes the
size as follows.
\begin{center}
\begin{tabular}{@{}lrrrr@{}}
\toprule
base & characteristic & Hilbert function & length & cardinality\\
\midrule
universal cubic base & \(4\) & \((1,15,107,509)\) & \(632\) & \(2^{632}\)\\
compressed base \(R\) & \(4\) & \((1,4,4,1)\) & \(10\) & \(2^{10}\)\\
\bottomrule
\end{tabular}
\end{center}

There are two different minimality statements to keep separate.
\begin{itemize}
\item The length-ten ring is optimal among \emph{all Artin quotients of the
audited universal cubic base} which retain either nonzero fourth-power
coordinate.  This is proved in Section~\ref{sec:minimality}.
\item No claim is made here that an unrelated rank-four counterexample over
an arbitrary base must have length at least ten.  Earlier exact work summarized
in \path{STRUCTURAL_LENGTH6_QUOTIENTS_F2_2026-07-10.md} (not re-audited in
this note) excludes residue-\(\F_2\) bases of length at most six; lengths
seven, eight, and nine remain the honest global frontier unless one proves
that every such counterexample is controlled by this cubic chart.
\end{itemize}

\section{The length-ten base ring}\label{sec:base}

\subsection{A compact four-generator presentation}

Define
\begin{equation}\label{eq:R-presentation}
R=(\Z/4)[x,y,z,d]/J,
\end{equation}
where
\begin{equation}\label{eq:J}
\begin{split}
J=\bigl(&2x,2y,2d,\ x^2-2,\ z^2-2,\ y^2,\ yz,\ zd,\\
        &d^2-xd,\ yd-xz-xy\bigr).
\end{split}
\end{equation}
Put
\begin{equation}\label{eq:secondary-names}
 \xi=xz,\qquad \eta=xd,\qquad \omega=xy,
 \qquad \sigma=2z.
\end{equation}
Every element has a unique additive normal form
\begin{equation}\label{eq:additive-normal-form}
 a+bz+\alpha x+\beta y+\gamma d
       +\lambda\xi+\mu\eta+\nu\omega,
\end{equation}
where \(a,b\in\Z/4\) and the six Greek coefficients lie in \(\F_2\).
Consequently
\begin{equation}\label{eq:additive-type}
 R_{\mathrm{add}}\simeq(\Z/4)^2\oplus(\Z/2)^6,
 \qquad |R|=2^{10}.
\end{equation}

\subsection{The inverse-system description}

The cleanest proof that the presentation is consistent and associative uses
a cubic inverse system.  Let
\[
 V=\F_2\{X,Y,Z,D\},\qquad V^\vee=\F_2\{X^\vee,Y^\vee,Z^\vee,D^\vee\},
\]
and let
\begin{equation}\label{eq:cubic}
 \mathcal G=X^{[2]}Z+XYD+YD^{[2]}+Z^{[3]}
\end{equation}
be a divided-power cubic.  Equivalently, its symmetric trilinear functional
is one precisely on the unordered triples
\[
 (X,X,Z),\quad(X,Y,D),\quad(Y,D,D),\quad(Z,Z,Z).
\]
Write
\[
 B_{\mathcal G}:\operatorname{Sym}^2(V)\longrightarrow V^\vee,
 \qquad B_{\mathcal G}(v,w)(t)=\mathcal G(v,w,t).
\]

Start with the additive group
\begin{equation}\label{eq:apolar-additive}
 \Z/4\{1\}\oplus\Z/4\{Z\}
 \oplus\F_2\{X,Y,D,X^\vee,Y^\vee,D^\vee\}.
\end{equation}
Identify
\begin{equation}\label{eq:ramified-identifications}
 Z^\vee=2,\qquad s=2Z.
\end{equation}
Give it a unital multiplication by declaring, for \(v,w\in V\) and
\(\lambda,\mu\in V^\vee\),
\begin{equation}\label{eq:apolar-products}
 vw=B_{\mathcal G}(v,w),\qquad
 v\lambda=\lambda(v)s,\qquad
 \lambda\mu=0,\qquad \m s=0.
\end{equation}
The identification \(Z^\vee=2\) is compatible with the additive group:
both rules give \(2v=0\) for \(v=X,Y,D\), and \(2Z=s\).
More explicitly, \(Z^\vee v=2v=Z^\vee(v)s\), while
\((2Z)v=2B_{\mathcal G}(Z,v)=0=sv\); the same compatibility against
\(V^\vee\) and \(s\) is immediate.  Hence the declared products descend
through both ramified identifications in \eqref{eq:ramified-identifications}.

Associativity is now transparent.  On three tangent vectors,
\[
 (v_1v_2)v_3=\mathcal G(v_1,v_2,v_3)s
             =v_1(v_2v_3)
\]
by symmetry of \(\mathcal G\).  If two factors lie in \(V^\vee\), or one
factor is \(s\), both bracketings vanish.  The remaining cases reduce to the
evaluation pairing in \eqref{eq:apolar-products}.  This constructs an
associative commutative ring of order \(2^{10}\).

Its nonzero products between tangent generators, up to symmetry, are
\begin{equation}\label{eq:tangent-products}
\begin{aligned}
 X^2&=Z^2=2,& XY&=D^\vee,& XZ&=X^\vee,& XD&=Y^\vee,\\
 YD&=X^\vee+D^\vee,& D^2&=Y^\vee.&&
\end{aligned}
\end{equation}
Upon putting \((X,Y,Z,D)=(x,y,z,d)\), this is exactly
\eqref{eq:R-presentation}--\eqref{eq:J}, with
\[
 X^\vee=\xi,\qquad Y^\vee=\eta,\qquad D^\vee=\omega.
\]
Conversely, the relations in \(J\) reduce every monomial to the span in
\eqref{eq:additive-normal-form}, so the presentation has at most \(2^{10}\)
elements.  Its map onto the just-constructed table ring has \(2^{10}\)
elements, proving both uniqueness of the normal form and the claimed
isomorphism.

\subsection{Locality, Hilbert function, and socle}

\begin{proposition}\label{prop:base-invariants}
The ring \(R\) is local Gorenstein of characteristic \(4\).  Its maximal
ideal and its powers are
\begin{align}
 \m&=(x,y,z,d),\nonumber\\
 \m^2&=\F_2\{2,\xi,\eta,\omega,\sigma\},\nonumber\\
 \m^3&=\F_2\{\sigma\},\qquad \m^4=0.\label{eq:m-powers}
\end{align}
Moreover
\[
 \Soc(R)=\F_2\sigma,\qquad H_R=(1,4,4,1),
 \qquad \length(R)=10.
\]
\end{proposition}

\begin{proof}
The quotient by \((x,y,z,d)\) is \(\F_2\).  In the normal form
\eqref{eq:additive-normal-form}, an element is a unit exactly when its
constant coefficient is odd, so this is the unique maximal ideal.  The
relations give \eqref{eq:m-powers} directly.

The degree-one and degree-two pairing into the top line is perfect.  In the
ordered bases \((x,y,z,d)\) and \((\xi,\eta,2,\omega)\), its diagonal
products are
\begin{equation}\label{eq:perfect-pairing}
 x\xi=y\eta=z\cdot2=d\omega=\sigma,
\end{equation}
and all off-diagonal products are zero.  Hence no nonzero class in
\(\m/\m^2\) or \(\m^2/\m^3\) annihilates \(\m\).  Since \(\m\sigma=0\),
the socle is exactly \(\F_2\sigma\).  The successive dimensions in
\eqref{eq:m-powers} are \((1,4,4,1)\), whose sum is ten.  Finally \(2\ne0\)
and \(4=0\), so the characteristic is exactly four.
\end{proof}

\section{The explicit rank-four Hopf algebra}\label{sec:hopf}

Let
\[
 A=R\{1,e_1,e_2,e_3\}
\]
be free of rank four.  Its augmentation sends \(e_1,e_2,e_3\) to zero.
Define a commutative multiplication by
\begin{equation}\label{eq:A-multiplication}
\begin{aligned}
 e_1^2&=xe_1+ye_2,& e_1e_2&=e_3,& e_1e_3&=xe_3,\\
 e_2^2&=ze_2,&e_2e_3&=ze_3,&e_3^2&=\xi e_3.
\end{aligned}
\end{equation}
Put
\begin{equation}\label{eq:q-sigma}
 q=x(y+z)=\xi+\omega,\qquad \sigma=2z.
\end{equation}
The coproduct is
\begin{equation}\label{eq:Delta1}
\begin{split}
\Delta(e_1)={}&e_1\otimes1+1\otimes e_1
 +x e_1\otimes e_1 +(y+z)e_2\otimes e_1+q e_3\otimes e_1,
\end{split}
\end{equation}
\begin{equation}\label{eq:Delta2}
\begin{split}
\Delta(e_2)={}&e_2\otimes1+1\otimes e_2
 +d e_2\otimes e_1 +(y+z)e_2\otimes e_2+q e_2\otimes e_3,
\end{split}
\end{equation}
and
\begin{equation}\label{eq:Delta3}
\begin{split}
\Delta(e_3)={}&e_3\otimes1+1\otimes e_3
 +e_1\otimes e_2+e_2\otimes e_1+x e_1\otimes e_3\\
& +(2+\eta)e_2\otimes e_1+q e_2\otimes e_2
  +\sigma e_2\otimes e_3\\
& +(x+d)e_3\otimes e_1+(y+z+\sigma)e_3\otimes e_2
  +q e_3\otimes e_3.
\end{split}
\end{equation}

\begin{proposition}\label{prop:hopf}
Equations \eqref{eq:A-multiplication}--\eqref{eq:Delta3}, together with
\(\Delta(1)=1\otimes1\), make \(A\) a counital bialgebra.  An antipode is
given by
\begin{equation}\label{eq:antipode}
\begin{aligned}
 S(e_1)&=e_1+\omega e_2+(y+z)e_3,\\
 S(e_2)&=e_2+d e_3,\\
 S(e_3)&=(3+\eta)e_3,\qquad S(1)=1.
\end{aligned}
\end{equation}
Consequently \(G=\Spec A\) is a finite locally free group scheme of rank
four over \(R\).
\end{proposition}

\begin{proof}[Exact verification]
All claims are identities in the 1024-element ring of composition length ten described in
Section~\ref{sec:base}.  The algebra associators reduce immediately from
\eqref{eq:A-multiplication}.  For example,
\[
 (e_1^2)e_2=xe_3+yz e_2=xe_3=e_1(e_1e_2),
\]
and
\[
 (e_1^2)e_3=x^2e_3+yz e_3=2e_3=e_1(e_1e_3).
\]
The remaining basis triples are of the same size.
The coordinate formula for multiplication in \(R\) is biadditive by
construction.  As an implementation gate, the verifier also checks the
distributive translation identity for every one of the \(1024\) ring
elements and every pair of additive generators.

For the coproduct, direct expansion verifies
\[
 (\Delta\otimes\id)\Delta(e_i)
 =(\id\otimes\Delta)\Delta(e_i),
 \qquad
 \Delta(e_ie_j)=\Delta(e_i)\Delta(e_j)
\]
for all basis indices.  The exact dependency-free verifier listed in
Section~\ref{sec:audit} checks all \(4^3=64\) algebra associators, all
\(4^2=16\) multiplicativity identities in the \(16\)-element tensor basis,
and all four coassociativity identities in the \(64\)-element triple-tensor
basis.  It then checks both convolution equations
\[
 \id*S=\iota\eps=S*\id
\]
for the explicit map \eqref{eq:antipode}.  The computation uses the literal
ring table and integer bit tuples; it invokes neither a solver nor a
Groebner basis.  A separate substitution into the independently exported
\(189\) universal bialgebra equations gives zero in every coordinate.
These finite coefficient checks prove the displayed identities exactly.
\end{proof}

Modulo \(\m\), the formulas reduce to
\[
 A/\m A\simeq\F_2[e_1,e_2]/(e_1^2,e_2^2),\qquad e_3=e_1e_2,
\]
with the usual coproduct of \(\alpha_2\times\alpha_2\).  Thus the failure
below is purely infinitesimal.

\section{Other descriptions of the group scheme}\label{sec:descriptions}

\subsection{A relative complete intersection and its functor of points}

The six products in \eqref{eq:A-multiplication} can be compressed to two
monic equations.

\begin{proposition}\label{prop:complete-intersection}
There is an augmented-algebra isomorphism
\begin{equation}\label{eq:CI-presentation}
 A\simeq R[U,V]/(U^2-xU-yV,\ V^2-zV),
 \qquad e_1=U,\quad e_2=V,\quad e_3=UV.
\end{equation}
In particular, the underlying finite flat scheme is an iterated pair of
monic quadratic covers, a relative complete intersection, and relative
Gorenstein.
\end{proposition}

\begin{proof}
The two monic relations reduce every polynomial uniquely to the
\(R\)-span of \(1,U,V,UV\).  They imply
\[
 U(UV)=xUV+yzV=xUV,\qquad V(UV)=zUV,
\]
and
\[
 (UV)^2=(xU+yV)(zV)=xz(UV),
\]
because \(yz=0\).  Thus they reproduce the full multiplication table
\eqref{eq:A-multiplication}; both sides are free of rank four, so the
displayed surjection is an isomorphism.  The final assertions are the
standard consequences of a monic codimension-two complete intersection.
\end{proof}

Consequently, for every commutative \(R\)-algebra \(T\),
\begin{equation}\label{eq:points}
 G(T)=\{(u,v)\in T^2:u^2=xu+yv,\ v^2=zv\}.
\end{equation}
Put \(w=uv\), and write a second point as
\(P'=(u',v',w')\), where \(w'=u'v'\).  Evaluating the coproduct gives
\begin{equation}\label{eq:point-law}
\begin{aligned}
 u''={}&u+u'+u'\bigl(xu+(y+z)v+q uv\bigr),\\
 v''={}&v+v'+v\bigl(du'+(y+z)v'+q u'v'\bigr),\\
 w''={}&u''v''.
\end{aligned}
\end{equation}
The identity is \((0,0)\), and the inverse is
\begin{equation}\label{eq:point-inverse}
 (u,v)^{-1}
 =\bigl(u+\omega v+(y+z)uv,\ v+duv\bigr).
\end{equation}
These formulas are polynomial identities subject to
\eqref{eq:points}; no choice of geometric points is involved.  Modulo
\(\m\), they become ordinary addition on
\(\alpha_2^2(T)=\{(u,v):u^2=v^2=0\}\).

\subsection{The dual measure Hopf algebra}

Let \(D=A^\vee=\operatorname{Hom}_R(A,R)\), let \(1_D=\eps\), and let
\(E_i\) be dual to \(e_i\) and vanish on \(1\).  Dualizing the coproduct
of \(A\) gives the following multiplication table:
\begin{equation}\label{eq:dual-multiplication}
\begin{array}{lll}
E_1^2=xE_1,&E_1E_2=E_3,
 &E_2E_1=(y+z)E_1+dE_2+(3+\eta)E_3,\\
E_1E_3=xE_3,&E_3E_1=qE_1+(x+d)E_3,
 &E_2^2=(y+z)E_2+qE_3,\\
E_2E_3=qE_2+\sigma E_3,
 &E_3E_2=(y+z+\sigma)E_3,&E_3^2=qE_3.
\end{array}
\end{equation}
Thus \(D\) is a PBW-type algebra with basis
\(1,E_1,E_2,E_1E_2\).  Its decisive reordering relation is
\begin{equation}\label{eq:dual-reordering}
 E_2E_1=(y+z)E_1+dE_2+(3+\eta)E_1E_2.
\end{equation}
In particular, noncommutativity can be read either from \(\Delta_A\) or
from multiplication in \(D\).

Dualizing the multiplication of \(A\) gives the cocommutative coalgebra
\begin{align}
 \Delta_D(E_1)={}&E_1\otimes1+1\otimes E_1+xE_1\otimes E_1,
 \label{eq:dual-Delta1}\\
 \Delta_D(E_2)={}&E_2\otimes1+1\otimes E_2
  +yE_1\otimes E_1+zE_2\otimes E_2,
 \label{eq:dual-Delta2}\\
 \Delta_D(E_3)={}&E_3\otimes1+1\otimes E_3
  +E_1\otimes E_2+E_2\otimes E_1\nonumber\\
 &+x(E_1\otimes E_3+E_3\otimes E_1)
  +z(E_2\otimes E_3+E_3\otimes E_2)
  +\xi E_3\otimes E_3.\label{eq:dual-Delta3}
\end{align}
The antipode is
\begin{equation}\label{eq:dual-antipode}
\begin{aligned}
 S_D(E_1)&=E_1,\\
 S_D(E_2)&=\omega E_1+E_2,\\
 S_D(E_3)&=(y+z)E_1+dE_2+(3+\eta)E_3.
\end{aligned}
\end{equation}
This is the finite distribution, or measure, Hopf algebra of \(G\).
Because its algebra is noncommutative, \(D\) is not the coordinate ring of
a Cartier-dual affine group scheme.  On the special fiber it becomes
\(\F_2[E_1,E_2]/(E_1^2,E_2^2)\), the usual measure algebra of
\(\alpha_2^2\).

\section{Faithful linear representations}\label{sec:linear}

For \(g=(u,v,w)\in G(T)\), right translation on functions is
\(r_g^*=(\id\otimes g)\Delta\).  In the basis
\((1,e_1,e_2,e_3)\), put
\begin{equation}\label{eq:KLM}
\begin{aligned}
 K&=1+du+(y+z)v+qw,\\
 L&=(3+\eta)u+qv+\sigma w,\\
 M&=1+(x+d)u+(y+z+\sigma)v+qw.
\end{aligned}
\end{equation}
Then
\begin{equation}\label{eq:GL4-matrix}
 \widetilde\rho(g)=
 \begin{pmatrix}
 1&u&v&w\\
 0&1+xu&0&v+xw\\
 0&(y+z)u&K&L\\
 0&qu&0&M
 \end{pmatrix}.
\end{equation}
Coassociativity gives
\(\widetilde\rho(gg')=\widetilde\rho(g)\widetilde\rho(g')\), and the
inverse matrix is \(\widetilde\rho(g^{-1})\).  Since the first row contains
the three universal coordinate functions, this already realizes \(G\) as a
closed subgroup of \(\GL_{4,R}\).

There is a smaller faithful representation.  The constant line \(R1\) is
fixed, so right translation descends to the free rank-three quotient
\begin{equation}\label{eq:reduced-regular}
 Q=A/R1=R\{\bar e_1,\bar e_2,\bar e_3\}.
\end{equation}
The resulting reduced right regular representation is
\begin{equation}\label{eq:GL3-matrix}
 \rho(g)=
 \begin{pmatrix}
 1+xu&0&v+xw\\
 (y+z)u&K&L\\
 qu&0&M
 \end{pmatrix}.
\end{equation}

\begin{theorem}\label{thm:GL3-minimal}
The morphism
\[
 \rho:G\longrightarrow\GL(Q)\simeq\GL_{3,R}
\]
is a closed immersion.  No closed immersion
\(G\hookrightarrow\GL_{2,R}\) exists.  Thus three is the minimal dimension
of a faithful finite locally free representation of \(G\).
\end{theorem}

\begin{proof}
Invertibility of \eqref{eq:GL3-matrix} follows either from
\eqref{eq:GL4-matrix}, which fixes \(R1\), or directly from
\(\rho(g^{-1})\).  To prove scheme-theoretic faithfulness, consider its two
matrix coefficients
\[
 \mathcal U=e_2+xe_3,\qquad
 \mathcal V=(3+\eta)e_1+qe_2+\sigma e_3.
\]
The base-ring identities
\[
 qx=\sigma,\qquad (3+\eta)^2=1
\]
give
\begin{equation}\label{eq:recover-e1}
 e_1=(3+\eta)(\mathcal V-q\mathcal U).
\end{equation}
Moreover
\[
 \mathcal U=(1+xe_1)e_2.
\]
If \(t=xe_1\), then \(t^4=0\) because \(x^4=0\).  Hence
\begin{equation}\label{eq:recover-e2}
 e_2=(1-t+t^2-t^3)\mathcal U,
 \qquad e_3=e_1e_2.
\end{equation}
Thus the matrix coefficients generate all of \(A\), which says precisely
that \(\mathcal O(\GL_3)\to A\) is surjective.

For minimality, base change to \(k=R/\m=\F_2\).  The representation becomes
\begin{equation}\label{eq:special-GL3}
 (u,v)\longmapsto
 \begin{pmatrix}1&0&v\\0&1&u\\0&0&1\end{pmatrix},
 \qquad G_k\simeq\alpha_2^2.
\end{equation}
A closed immersion \(\alpha_2^2\hookrightarrow\GL_{2,k}\) would inject
restricted Lie algebras and therefore give two linearly independent,
commuting, square-zero matrices in \(M_2(k)\).  A nonzero square-zero
matrix is conjugate to \(E_{12}\).  Its centralizer consists of
\(aI+bE_{12}\), and the square-zero elements in this centralizer are only
the multiples \(bE_{12}\).  Such an independent pair cannot exist.
\end{proof}

The representation itself has a useful filtration:
\begin{equation}\label{eq:module-filtration}
 0\subset R\bar e_2\subset Q.
\end{equation}
Indeed, the middle column of \eqref{eq:GL3-matrix} is
\((0,K,0)^{\mathsf t}\).  In the reordered basis
\((\bar e_2,\bar e_1,\bar e_3)\), the image lies in the parabolic stabilizing
a line.  This is a filtration of a \(G\)-module, not a filtration of \(G\)
by subgroup schemes.  On the special fiber, the invariant line is trivial
and sees neither \(\alpha_2\)-direction; the two-dimensional quotient sees
only the \(v\)-direction.  The off-diagonal extension entry
\(L\bmod\m=u\) records the other direction.

\section{Rank-two subgroups, cosets, and filtrations}
\label{sec:subgroups}

\subsection{An explicit Oort--Tate subgroup and its coset space}

Setting \(e_2=e_3=0\) in the Hopf algebra gives a surjection
\begin{equation}\label{eq:basic-H-quotient}
 A\longrightarrow B_x=R[T]/(T^2-xT),\qquad
 e_1\mapsto T,\quad e_2,e_3\mapsto0,
\end{equation}
where
\begin{equation}\label{eq:basic-H-Hopf}
 \Delta(T)=T\otimes1+1\otimes T+xT\otimes T,
 \qquad S(T)=T.
\end{equation}
Thus it defines a rank-two closed subgroup \(H_x\subset G\).  In the
normalization used here, this is the self-dual Oort--Tate group with
parameter pair \((x,x)\); the parameter relation is \(x^2=2=-2\).  It is
killed by two because
\begin{equation}\label{eq:H-killed-two}
 [2]^\#(T)=2T+xT^2=(2+x^2)T=0.
\end{equation}

The subgroup has a concrete quotient as a homogeneous space.  Put
\begin{equation}\label{eq:left-coset-algebra}
 C=R[V]/(V^2-zV).
\end{equation}
By \eqref{eq:CI-presentation}, \(A=C[U]/(U^2-xU-yV)\), free of rank two
over \(C\).  The left \(H_x\)-coaction is
\begin{equation}\label{eq:left-H-coaction}
 \lambda(V)=1\otimes V,\qquad
 \lambda(U)=T\otimes1+1\otimes U+xT\otimes U.
\end{equation}
It follows that \(C=A^{H_x}\), and
\begin{equation}\label{eq:left-coset}
 H_x\backslash G\simeq\Spec C.
\end{equation}
For completeness, the torsor identity can be seen without invoking a
general quotient theorem.  Since \((1+xU)^2=1\),
\[
 T\otimes1=(\lambda(U)-1\otimes U)(1\otimes(1+xU)).
\]
This makes the canonical torsor map surjective; its source and target are
finite locally free of the same rank, hence it is an isomorphism.

Although \(C\) has the shape of an Oort--Tate algebra, it is not a quotient
\emph{group} algebra of \(G\).  Indeed,
\begin{equation}\label{eq:e2-external-terms}
 \Delta(e_2)=e_2\otimes1+1\otimes e_2
 +d e_2\otimes e_1+(y+z)e_2\otimes e_2+q e_2\otimes e_3
\end{equation}
has terms outside \(C\otimes C\).  Thus \(H_x\subset G\) is a
nonnormal Oort--Tate torsor fibration, not an extension of two rank-two
groups.

\subsection{Classification of all rank-two subgroups over
\texorpdfstring{\(R\)}{R}}

There are several embeddings of the same abstract subgroup.  For
\((\epsilon,\alpha,\beta,\gamma)\in\F_2^4\), put
\begin{equation}\label{eq:subgroup-parameters}
\begin{aligned}
 r&=\epsilon d+\alpha\eta+\beta q+\gamma\sigma,\\
 t&=x+\epsilon q+\alpha\sigma,\\
 s&=\epsilon\eta+(\epsilon+\beta)\sigma.
\end{aligned}
\end{equation}
Here the coefficient \(\epsilon+\beta\) is taken in \(\F_2\).

\begin{theorem}\label{thm:subgroup-classification}
For every quadruple in \eqref{eq:subgroup-parameters}, \(t^2=2\), and
\[
 B_t=R[f]/(f^2-tf),\qquad
 \Delta(f)=f\otimes1+1\otimes f+t f\otimes f,\qquad S(f)=f
\]
is a rank-two Hopf algebra.  The assignments
\begin{equation}\label{eq:all-subgroup-maps}
 e_1\mapsto f,\qquad e_2\mapsto rf,\qquad e_3\mapsto sf
\end{equation}
are Hopf surjections with pairwise distinct kernels
\begin{equation}\label{eq:all-subgroup-ideals}
 I_{\epsilon,\alpha,\beta,\gamma}
 =(e_2-re_1,\ e_3-se_1).
\end{equation}
These are exactly the rank-two finite locally free closed subgroups of
\(G\) defined over \(R\).  All sixteen are abstractly isomorphic to
\(H_x\), and all reduce to the same \(\alpha_2\)-line
\(e_2=e_3=0\) in \(G_k=\alpha_2^2\).
\end{theorem}

\begin{proof}
Direct substitution of \eqref{eq:subgroup-parameters} into
\eqref{eq:A-multiplication}--\eqref{eq:Delta3} proves that
\eqref{eq:all-subgroup-maps} is a Hopf map.  Here is also a short
completeness argument.  For any rank-two Hopf quotient, the augmentation
splitting and locality of \(R\) give
\[
 B=R\oplus Rf,\qquad f^2=tf,\qquad
 \Delta(f)=f\otimes1+1\otimes f+c f\otimes f.
\]
At least one of the images of \(e_1,e_2\) generates \(Rf\).  The branch in
which only the image of \(e_2\) is a generator is impossible modulo
\(\m^2\): after normalizing \(e_2\mapsto f\), one has \(t=z\), and
\(e_1\mapsto af\), with \(a\in\m\), would force
\(a^2z=xa+y\), whose two sides have incompatible degree-one parts.

Normalize therefore \(e_1\mapsto f\), and write \(e_2\mapsto rf\).  The
algebra equations first give
\begin{equation}\label{eq:subgroup-derivation-one}
 t=x+yr,\qquad e_3\mapsto rt\,f.
\end{equation}
The relation \(e_2^2=ze_2\) then forces, in the additive normal basis,
\[
 r=\lambda y+\epsilon d+\rho\xi+\alpha\eta
       +\tau\omega+\gamma\sigma.
\]
Comparing the \(f\otimes f\) coefficient in \(\Delta(e_2)\) forces
\(\lambda=0\) and \(\rho=\tau\).  Thus \(r\) is exactly the first line of
\eqref{eq:subgroup-parameters}; \eqref{eq:subgroup-derivation-one} and the
remaining equations give the displayed \(t,s\), and also \(c=t\).

Finally, put
\[
 a_0=1+\epsilon(y+z)+\alpha\xi.
\]
Then \(t=a_0x\) and \(a_0^2x=x\).  The change of generator
\(f\mapsto a_0T\) identifies \(B_t\) with \(B_x\), including coproducts.
Distinct pairs \((r,s)\) give distinct kernels, so there are exactly
sixteen.
\end{proof}

This classification is over the original base \(R\).  After a
non-faithfully-flat specialization the group can acquire additional
subgroups: for example, \(G_k=\alpha_2^2\) has the usual projective line of
\(\alpha_2\)-subgroups.  Over \(\F_2\), its other two rational lines do not
lift to \(R\); the one line which does lift has the sixteen infinitesimally
distinct embeddings above.

\subsection{Why none is normal, and why no composition filtration exists}

There is a direct certificate for the failure of normality.  If
\(q_r:A\to B_t\) denotes \eqref{eq:all-subgroup-maps} and
\[
 \Ad^*(a)=\sum a_{(2)}\otimes a_{(1)}S(a_{(3)}),
\]
then exact expansion gives
\begin{equation}\label{eq:all-nonnormal-certificate}
 (q_r\otimes\id)\Ad^*(e_3-se_1)
 =f\otimes(\eta+\epsilon\sigma)e_2\ne0.
\end{equation}
Thus none of the sixteen subgroups is normal.  There is also a more
structural argument, which remains valid after faithfully flat base change.

\begin{lemma}\label{lem:rank-two-killed-two}
Every finite locally free group scheme of rank two is killed by two.
\end{lemma}

\begin{proof}
The assertion is local on the base.  In a rank-two augmented Hopf algebra,
choose a basis \(1,T\) with
\[
 T^2=aT,\qquad
 \Delta(T)=T\otimes1+1\otimes T+bT\otimes T.
\]
Multiplicativity of \(\Delta\) gives
\((ab+1)(ab+2)=0\).  If \(S(T)=cT\), the antipode equation says
\(1+c(1+ab)=0\), so \(1+ab\) is a unit.  Hence \(ab=-2\), and
\[
 [2]^\#(T)=2T+bT^2=(2+ab)T=0.
\]
\end{proof}

If a rank-two subgroup \(H\subset G\) were normal, the fppf quotient
\(G/H\) would also have rank two.  Both \(H\) and \(G/H\) would be killed
by two.  The image of any section \(g\) in the quotient would satisfy
\(\bar g^2=1\), so fppf-locally \(g^2\in H\), and hence \(g^4=1\).
This would make \([4]\) trivial on \(G\), contradicting
\eqref{eq:fourth-power}.  Faithful flatness preserves the nonzero free-basis
coefficient \(\sigma\), so a normal rank-two subgroup cannot appear even
after a faithfully flat extension of \(R\).

Finally, the ranks in a finite locally free normal filtration multiply.
A nontrivial proper step in rank four would therefore have rank two.  We
conclude:
\begin{center}
\emph{The group \(G\) has rank-two subgroups, but no nontrivial normal or}\\
\emph{composition filtration by finite flat group schemes.}
\end{center}
The chains \(1\subset H\subset G\) exist as subgroup chains, and
\(H\backslash G\) is a rank-two homogeneous space, but there is no
rank-two quotient group.

\section{The power words by hand}\label{sec:powers}

Let \(m_A:A\otimes_R A\to A\) be multiplication and set
\[
 \phi=m_A\circ\Delta=[2]^\#.
\]
Contracting \eqref{eq:Delta1}--\eqref{eq:Delta3} with the multiplication
table gives the remarkably small formulas
\begin{equation}\label{eq:phi}
\begin{aligned}
 \phi(e_1)&=\omega e_2+(y+z+\sigma)e_3,\\
 \phi(e_2)&=d e_3,\\
 \phi(e_3)&=0.
\end{aligned}
\end{equation}
For example, the two primitive legs and the term \(x e_1\otimes e_1\)
contribute
\[
 2e_1+x e_1^2=2e_1+(2e_1+\omega e_2)=\omega e_2,
\]
while
\[
 q(e_3e_1)=qx e_3=x^2(y+z)e_3=\sigma e_3.
\]
This proves the first line of \eqref{eq:phi}; the other two follow by the
same direct contraction.

Since \([4]^\#=\phi\circ\phi\), one obtains
\begin{equation}\label{eq:fourth-power}
\begin{aligned}
 [4]^\#(e_1)&=\omega d\,e_3=xyd\,e_3=\sigma e_3,\\
 [4]^\#(e_2)&=0,\qquad [4]^\#(e_3)=0.
\end{aligned}
\end{equation}
Indeed, \(yd=xz+xy\) gives
\begin{equation}\label{eq:socle-carry}
 xyd=x(xz+xy)=x^2z+x^2y=2z+2y=\sigma.
\end{equation}
The element \(\sigma=2z\) is the nonzero socle generator, so
\eqref{eq:fourth-power} proves that \(G\) is not killed by four.

There is also an exact quasi-idempotence relation.  Let
\(\iota:R\to A\) denote the unit and let
\[
 P=\iota\eps,\qquad \psi=\phi-P.
\]
On the augmentation ideal, \(P=0\), and \eqref{eq:phi} gives
\begin{equation}\label{eq:psi-relation-small}
 \psi^2=2\psi.
\end{equation}
Indeed, \(2\omega=2d=0\) and
\(2(y+z+\sigma)=2z=\sigma\).  It follows that
\[
 \psi^3=4\psi=0.
\]
Moreover \(P^2=P\) and \(P\psi=\psi P=0\).  Therefore
\(\phi^3=(P+\psi)^3=P+\psi^3=P\), which is equivalently
\begin{equation}\label{eq:eighth}
 [8]^\#=P.
\end{equation}
Thus the exact pointwise power-word exponent is eight.

Finally, the group is genuinely noncommutative.  In \(\Delta(e_1)\), the
coefficient of \(e_1\otimes e_2\) is zero, whereas the coefficient of
\(e_2\otimes e_1\) is \(y+z\ne0\).  Hence \(\Delta\) is not
cocommutative.

\section{Conceptual explanation: carry times cocommutator}
\label{sec:conceptual}

\subsection{The intrinsic first symbol}

Let \(A_0=A/\m A\), let \(J_0=\ker(\eps:A_0\to\F_2)\), and put
\[
 V_0=J_0/J_0^2,\qquad L=J_0^2=\F_2e_3.
\]
The first \(\m\)-adic symbol of \(\psi\) is
\[
 N\in\gr_1(R)\otimes_{\F_2}\operatorname{Hom}_{\F_2}(V_0,L)
\]
with
\begin{equation}\label{eq:N}
 N(\bar e_1)=(\bar y+\bar z)\otimes e_3,\qquad
 N(\bar e_2)=\bar d\otimes e_3.
\end{equation}
Here bars on coefficients denote classes in \(\gr_1(R)\).  Extend \(N\) by
zero on \(L\subset J_0\); then it is an element of
\(\gr_1(R)\otimes\operatorname{End}_{\F_2}(J_0)\), and its graded square is
zero because its image lies in \(L\).  This is the conceptual reason that
the putative quadratic fourth-power obstruction cancels.

The mixed-characteristic secondary class is
\[
 \delta=\operatorname{in}_2(2)\in\m^2/\m^3.
\]
The first nonzero fourth-power symbol is
\begin{equation}\label{eq:deltaN}
 \operatorname{in}_3(\psi^2)=\delta N
 \in\gr_3(R)\otimes_{\F_2}\operatorname{Hom}_{\F_2}(V_0,L).
\end{equation}
For this example, \(2=x^2=z^2\), and
\[
 (\delta N)(\bar e_1)=\overline{2(y+z)}\otimes e_3
   =\bar\sigma\otimes e_3,\qquad
 (\delta N)(\bar e_2)=\overline{2d}\otimes e_3=0.
\]
Unlike a particular monomial normal form, the factors \(\delta\), \(N\),
and their product are invariant under a change of frame.

\subsection{The Oort--Tate carry and the skew direction}

In the universal coordinates, the decisive normal class was
\begin{equation}\label{eq:universal-normal-class}
 m_{11}^{\,1}c_{111}(c_{112}+c_{121}).
\end{equation}
In the compressed ring,
\begin{equation}\label{eq:mechanism-one-line}
 m_{11}^{\,1}=c_{111}=x,\qquad
 c_{112}+c_{121}=y+z,
\end{equation}
and therefore
\[
 m_{11}^{\,1}c_{111}(c_{112}+c_{121})
 =x^2(y+z)=2(y+z)=2z=\sigma.
\]

The product \(m_{11}^{\,1}c_{111}=2\) is the characteristic-two
binomial carry appearing in the compatibility equation
\(\Delta(e_1^2)=\Delta(e_1)^2\).  It is the same local mechanism as the
familiar multiplication--coproduct parameter pairing in a rank-two
Oort--Tate chart.  The second factor is a first-order skew coproduct
coefficient: in characteristic two, the sum and difference of
\(c_{112}\) and \(c_{121}\) have the same initial class.  Thus mixed
characteristic supplies the carry, while noncommutativity supplies the
direction on which it acts.

This also explains two necessary features of the example.
\begin{itemize}
\item Setting \(2=0\) kills \(\delta\) and hence the defect.
\item Imposing cocommutativity kills the skew map \(N\) and hence the defect.
\end{itemize}
The construction is therefore neither an equal-characteristic accident nor
a hidden commutative example.

\subsection{The exact universal identity}

The relation \eqref{eq:psi-relation-small} is not a coincidence of the
chosen quotient.  Exact filtered reduction of all nine coordinates on the
universal cubic base proves
\begin{equation}\label{eq:universal-psi}
 \boxed{\ \psi^2=2\psi\ }.
\end{equation}
Equivalently,
\[
 [4]^\#-P=2([2]^\#-P).
\]
The supplied checker reduces every coordinate of
\(T_{ir}-2\phi_i^{\,r}\) through the complete degree-three filtered ideal;
all nine cubic remainders are zero.  Formula \eqref{eq:deltaN} is the
associated-graded shadow of this exact identity.

\subsection{A deformation tower: normal extension, nonnormal torsor,
and fourth-power defect}\label{sec:deformation-tower}

The construction becomes more transparent when reduced along two different
filtrations of the base.  The \(\m\)-adic filtration records the order at
which the power obstruction appears.  Modulo \(\m\), the group is the
commutative group \(\alpha_2^2\).  Modulo \(\m^2\), its reduced doubling map
is already nonzero:
\begin{equation}\label{eq:first-order-doubling}
 \psi(e_1)=(y+z)e_3,\qquad
 \psi(e_2)=de_3,\qquad \psi(e_3)=0
 \pmod{\m^2}.
\end{equation}
This is precisely the symbol \(N\) of \eqref{eq:N}.  Its image lies in the
special-fiber square \(L=\F_2e_3\), on which it vanishes, so its first-order
square is zero.  The group is therefore still killed by four at this order.

The same tensor is the first skew coproduct symbol.  If
\(a,b\in\F_2\), then
\begin{equation}\label{eq:first-skew-symbol}
\begin{split}
 (\Delta-\tau\Delta)(ae_1+be_2)
 \equiv{}&\bigl(a(\bar y+\bar z)+b\bar d\bigr)\\
 &\otimes(e_2\otimes e_1-e_1\otimes e_2)
 \pmod{\m^2}.
\end{split}
\end{equation}
Since \(\bar y+\bar z\) and \(\bar d\) are independent in
\(\m/\m^2\), this map has no nonzero tangent kernel.  Thus the first-order
deformation has already destroyed every possible cocommutative rank-two
quotient direction, even though it has not yet produced a fourth-power
defect.  To spell out the quotient argument, a normal rank-two subgroup
would give a rank-two quotient \(K\), hence a rank-two Hopf subalgebra
\(\mathcal O(K)\subset A\).  Its special-fiber augmentation generator must
have a nonzero tangent term \(ae_1+be_2\): a pure \(e_3\)-term cannot
generate a Hopf subalgebra because \(\Delta_0(e_3)\) contains
\(e_1\otimes e_2+e_2\otimes e_1\).  A rank-two Hopf algebra has
cocommutative coproduct, while corrections in \(\m A\) have cocommutative
special-fiber coproduct and cannot cancel \eqref{eq:first-skew-symbol}.
The injectivity just proved is therefore a contradiction.  Since passing to
\(R_8\) does not change \(\m/\m^2\), this also excludes every normal
rank-two subgroup over the intermediate length-eight base, not merely the
particular subgroup \(H_x\).

Killing the two skew directions gives the length-four quotient
\begin{equation}\label{eq:commutative-quotient}
 R/(d,y+z)\simeq\F_2[x,z]/(x^2,z^2),
\end{equation}
over which the entire coproduct becomes cocommutative.  This is one precise
sense in which the original group is a noncommutative deformation of a
simpler commutative one.

At the opposite end, \(\m^3=\F_2\sigma\).  Hence
\begin{equation}\label{eq:socle-reduction}
 \bar R=R/(\sigma)=R/\m^3
\end{equation}
has length nine, and \(G_{\bar R}\) is killed by four.  The surjection
\(R\to\bar R\) has square-zero kernel.  Lifting across this final socle
extension changes the fourth-power map from zero to
\(\sigma e_3\).  Thus the counterexample itself is a first-order
deformation, in the square-zero direction \(\sigma\), of a rank-four group
which is killed by four.  The exact identity \(\psi^2=2\psi\) says that the
obstruction to lifting the relation \([4]=e\) is the Bockstein-type carry
applied to the first-order skew class; no separate connecting morphism is
being asserted by this terminology.

There is a second tower which isolates normality even more sharply.  Let
\begin{equation}\label{eq:d-ideal}
 I=(d)=\F_2\{d,\eta,q,\sigma\}.
\end{equation}
The base relations give
\begin{equation}\label{eq:d-ideal-powers}
 I^2=\F_2\{\eta,\sigma\},\qquad I^3=0,
 \qquad d^2=\eta,\qquad dq=\omega d=\sigma.
\end{equation}
Consequently the quotients
\begin{equation}\label{eq:6-8-10-tower}
 R_{10}=R\twoheadrightarrow
 R_8=R/I^2\twoheadrightarrow R_6=R/I
\end{equation}
have the indicated lengths, and both successive kernels are square-zero.

Over the length-six base, the group has the simpler form anticipated in the
question.  Explicitly,
\begin{equation}\label{eq:R6-presentation}
 R_6\simeq
 (\Z/4)[x,y,z]/
 (2x,2y,2z,x^2-2,z^2-2,y^2,yz,x(y+z)).
\end{equation}
Here \(d=\eta=q=\sigma=0\), so \(R_6\{1,e_2\}\subset A_{R_6}\) is a Hopf
subalgebra:
\begin{equation}\label{eq:K6}
\begin{aligned}
 e_2^2&=ze_2,\\
 \Delta(e_2)&=e_2\otimes1+1\otimes e_2
 +(y+z)e_2\otimes e_2,\\
 S(e_2)&=e_2.
\end{aligned}
\end{equation}
It is the coordinate algebra of a rank-two Oort--Tate quotient \(K_6\),
with parameter pair \((z,y+z)\); indeed \(z(y+z)=2\).  There is an exact
sequence of finite flat group schemes
\begin{equation}\label{eq:R6-extension}
 1\longrightarrow (H_x)_{R_6}
 \longrightarrow G_{R_6}\longrightarrow K_6\longrightarrow1.
\end{equation}
Both outer terms are killed by two, so \(G_{R_6}\) is killed by four.
Thus the length-six reduction really does have a rank-two-by-rank-two
normal filtration.

The first square-zero lift, from \(R_6\) to \(R_8\), turns on the classes
\(d\) and \(q=yd\).  The external terms
\begin{equation}\label{eq:normality-breaking-terms}
 d e_2\otimes e_1+q e_2\otimes e_3
\end{equation}
in \(\Delta(e_2)\) mean that \(R_8\{1,e_2\}\) is no longer a Hopf
subalgebra; the quotient group in \eqref{eq:R6-extension} has deformed into
only a coset space.  Dually, the subgroup \(H_x\) corresponds to the Hopf
subalgebra \(R\{1,E_1\}\subset D\).  Directly from
\eqref{eq:dual-multiplication}--\eqref{eq:dual-Delta3}, its adjoint
obstructions are
\begin{equation}\label{eq:dual-adjoint-normality}
\begin{aligned}
 \operatorname{ad}_{E_2}(E_1)
  &=(y+z)E_1+dE_2+\eta E_3,\\
 \operatorname{ad}_{E_3}(E_1)
  &=qE_1+\eta E_2.
\end{aligned}
\end{equation}
They lie in \(R\{1,E_1\}\) after a base change precisely when \(d=0\).
This proves that \(d\) is the deformation parameter which breaks
normality.  Since \(\sigma=0\) on \(R_8\), however, this intermediate group
is still killed by four.

The last square-zero lift, from \(R_8\) to \(R_{10}\), turns on
\(\eta=d^2\) and \(\sigma=dq=\omega d\).  Formula \eqref{eq:phi} now has
the composable path
\begin{equation}\label{eq:deformation-path}
 e_1\ \xmapsto{\ \omega\ }\ e_2
 \ \xmapsto{\ d\ }\ e_3,
 \qquad \omega d=\sigma,
\end{equation}
and hence \(\phi^2(e_1)=\sigma e_3\).  In summary,
\begin{center}
\begin{tabular}{@{}llll@{}}
\toprule
base & length & subgroup structure & power relation\\
\midrule
\(R_6=R/(d)\) & \(6\) & rank-two-by-rank-two extension & \([4]=e\)\\
\(R_8=R/(\eta,\sigma)\) & \(8\) & normality broken by \(d\) & \([4]=e\)\\
\(R\) & \(10\) & only nonnormal rank-two subgroups & \([4]\ne e\), \([8]=e\)\\
\bottomrule
\end{tabular}
\end{center}
The deformation therefore breaks the normal filtration one infinitesimal
layer before it breaks the fourth-power relation.  This separation is the
most concrete deformation-theoretic explanation of why the final group is
both non-filterable and not killed by its rank.

\section{How the smaller ring was found}\label{sec:discovery}

\subsection{From a cubic defect to an inverse system}

Let \(B\) be the previously audited universal cubic base.  Its maximal ideal
satisfies \(\m_B^4=0\), and
\[
 H_B=(1,15,107,509).
\]
The two possible nonzero fourth-power coordinates lie in
\(\m_B^3\).  To retain one while shrinking the base, choose an additive
functional on \(\m_B^3\) which is nonzero on that coordinate and annihilates
the cubic ideal relations.  Passing to its annihilator produces a
Gorenstein quotient.  On the associated graded, this functional is a cubic
inverse system; the rank of its first catalecticant determines the sizes of
the degree-one and degree-two layers.

The first \(38\)-term dual certificate from the universal audit had
catalecticant rank seven and led to a valid quotient of Hilbert function
\((1,7,7,1)\), hence length sixteen.  Varying the dual functional, rather
than merely quotienting that first ring, lowers the rank to four.  A
particularly sparse representative is the cubic \(\mathcal G\) in
\eqref{eq:cubic}.

Its four first contractions are
\begin{equation}\label{eq:contractions}
 XZ+YD,\quad XD+D^{[2]},\quad
 X^{[2]}+Z^{[2]},\quad XY+YD.
\end{equation}
They are linearly independent, so the first catalecticant has rank four.
The corresponding apolar Hilbert function is therefore
\[
 (1,4,4,1).
\]
The filtered identifications \(Z^\vee=2\) and \(s=2Z\) lift this homogeneous
apolar algebra to the characteristic-four ring in Section~\ref{sec:base}.

\subsection{Exact specialization of the universal parameters}

Appendix~\ref{app:specialization} gives the complete \(45\)-parameter map.
Substitution into the raw integral export makes every one of the \(189\)
bialgebra equations vanish and gives the target vector
\[
 (0,0,\sigma,0,0,0,0,0,0).
\]
This verifies that the inverse-system construction lifts exactly; it is not
merely a tangent-cone or associated-graded solution.
Every displayed parameter lies in \(\m_R\), so the fourth power of the
universal augmentation ideal maps to zero by \(\m_R^4=0\); the \(189\)
vanishing equations kill the universal bialgebra ideal.  Hence the
specialization factors through a map \(B\to R\).  It is surjective because
\(p_0,p_1,p_{10},p_{30}\) map respectively to the ring generators
\(x,y,z,d\).  Thus \(R\) is genuinely an Artin quotient of \(B\).

\subsection{Search chronology and methodological lessons}
\label{sec:search-process}

The construction emerged from an attempted proof of the killed-by-rank
assertion, and the reversal in viewpoint is worth recording.  The initial
program studied universal deformation charts around the local--local
rank-four fibers and expanded the reduced doubling operator in the
maximal-ideal filtration.  Its first symbol \(N\) always has square zero.
This explains the persistent cancellation of the naively quadratic
fourth-power obstruction, but it does not imply that the complete filtered
fourth-power class vanishes.  Considerable equal-characteristic and
low-depth work proved genuine special cases; resource failures and vanishing
in those cases were never a proof of the mixed-characteristic statement.

The decisive change was to preserve the integral filtration.  If
\(\mathfrak q=(2,p_0,\ldots,p_{44})\) in the universal chart, the relevant order of
an integral term is
\[
 \operatorname{ord}_{\mathfrak q}(c\,p^\alpha)=v_2(c)+|\alpha|.
\]
Integer terms must be added before this initial order is taken.  In
particular, two equal odd terms produce \(2f\), one filtration step higher;
reducing modulo two first would replace \(f+f\) by zero and erase precisely
the carry which produces the example.  With this convention, the quadratic
targets still cancel, but two cubic fourth-power targets on the
\(\alpha_2^2\)-chart survive.  At that point the project changed from
trying to prove universal membership to constructing exact nonmembership
certificates.

A nonzero normal form for one elimination order was not considered adequate
evidence.  The first audit instead produced an \(\F_2\)-linear functional
which annihilated all \(45{,}414\) exhaustive cubic ideal candidates and
evaluated to one on the target.  A second program parsed the integral
Macaulay2 export, used a different polynomial representation and elimination
implementation, and reconstructed independent dual certificates.  A
separate comparison checked the integral coefficients of all \(189\)
equations and nine targets before taking initial forms.  Only after these
data paths agreed was the counterexample interpretation accepted.

The first target-detecting \(38\)-term functional had catalecticant rank
seven and therefore produced the length-sixteen Gorenstein quotient
mentioned above.  That witness was not canonical: sparse support and small
catalecticant rank are different objectives.  Moreover, every proper
quotient of a fixed Artin Gorenstein witness kills its socle target, so
trying to compress that one ring further was the wrong search space.  The
successful move was to vary the functional in the full space of cubics
annihilating the universal relations and detecting a target.  Minimizing
the first catalecticant in that space produced the rank-four cubic
\(\mathcal G\).  Its apolar algebra suggested the small ring, while the
exact \(45\)-parameter substitution proved that the result was a genuine
inhomogeneous Hopf-algebra specialization rather than only an
associated-graded model.

Once the compact table had been found, the proof was deliberately rebuilt
without the universal elimination.  The \(1024\)-element tuple verifier
checks the ring and Hopf identities directly, and a second additive
coordinate model independently checks the same theorem.  The universal
chart now explains provenance, and the MinRank calculation addresses
relative compression; neither is needed for the finite direct proof that
\([4]^\#(e_1)=\sigma e_3\ne0\).

Finally, the conceptual descriptions were obtained by ablation after the
example was secure.  Killing the first skew directions makes the group
commutative.  Killing the socle class \(\sigma\) restores \([4]=e\).
Killing \(d\) restores a normal Oort--Tate filtration.  These reductions led
to the tower in Section~\ref{sec:deformation-tower} and separate the
first-order obstruction to normality from the later power-word obstruction.

The main methodological lessons are:
\begin{itemize}
\item preserve mixed-characteristic carries before passing to an associated
  graded ring;
\item optimize an invariant such as catalecticant rank, rather than the
  superficial number of variables in one presentation;
\item use associated-graded calculations for discovery, but lift every
  candidate through the exact inhomogeneous equations;
\item replace a large universal proof by finite direct arithmetic once a
  small model is known, and use genuinely independent implementations;
\item classify timeouts, out-of-memory exits, missing output, and solver
  \texttt{unknown} as inconclusive rather than as mathematical evidence;
\item separate unconditional existence from solver-dependent relative
  minimality.
\end{itemize}

These qualifications also delimit the next search.  Conditional on the
recorded MinRank verdict, no shorter quotient exists inside this cubic
chart.  A global example of length seven, eight, or nine would have to use a
different admissible dual, greater Loewy depth, another special fiber or
chart, or a mechanism not visible at cubic order.

\section{Optimality within the universal cubic construction}
\label{sec:minimality}

\subsection{The global MinRank calculation}

Let
\[
 V_B=\gr_1(B),\qquad W_B=\gr_2(B),\qquad U_B=\gr_3(B).
\]
Their dimensions are \((15,107,509)\).  For a cubic functional
\(L\in U_B^\vee\), form the first catalecticant
\[
 C_L:V_B\longrightarrow W_B^\vee.
\]
The exact question is whether some \(L\), nonzero on either fourth-power
target, has \(\operatorname{rank}C_L\le3\).

The computation reconstructs the \(15\) tangent generators and the \(171\)
independent relations in the \(680\)-dimensional space
\(\operatorname{Sym}^3(\F_2^{15})\).  In the ambient \(120\)-dimensional
quadratic monomial space, rank at most three is encoded without loss by
\begin{equation}\label{eq:rank-factorization}
 C=U_{15\times3}V_{3\times120}.
\end{equation}
Symmetry/integrability, annihilation of all \(171\) cubic relations, and
nonvanishing on at least one of the two targets are then Boolean equations.
Every matrix of rank at most three admits \eqref{eq:rank-factorization}, so
the encoding is complete rather than heuristic.

The resulting instance has \(405\) factor variables and \(1292\) assertions.
Z3~4.16.0 returns \texttt{UNSAT}.  The audited run used a 120-second timeout
and a 1024-MiB memory ceiling and completed in \(18.64\) seconds.  The
canonical SMT payload hash is
\[
\texttt{7f20f68bb134f26bbbe943089f0a92fcec206a0ebb48da48db528f3f76a4f1df}.
\]
This hash makes the encoding reproducible; it is not a separately checked
solver proof object.  Accordingly, the minimality conclusion below is a
computer-assisted conclusion conditional on the correctness of the Z3
UNSAT verdict.
The cubic \(\mathcal G\) supplies a rank-four witness.  Thus four is the
minimum catalecticant rank among \emph{all} cubic duals of this universal
chart which detect either target.

\subsection{From catalecticant rank to quotient length}

\begin{theorem}[Computer-assisted]\label{thm:minimal-quotient}
If \(D\) is an Artin quotient of the universal cubic base \(B\) and a
fourth-power target remains nonzero in \(D\), then
\[
 \length(D)\ge10.
\]
The ring \(R\) constructed above attains equality.
\end{theorem}

\begin{proof}
Quotient \(D\) further by an ideal maximal among those which do not kill the
chosen target \(\theta\).  In the resulting ring \(D'\), every nonzero ideal
contains \(\theta\).  Hence the socle is the one-dimensional line generated
by \(\theta\), so \(D'\) is Artin Gorenstein.  More explicitly,
\(\theta\in\n^3\) and \(\n^4=0\), hence \(\n\theta=0\); maximality then
forces every nonzero socle element to span the same line.
Since \(B\) has fourth power of its maximal ideal zero and \(\theta\) is a
nonzero cubic class,
\[
 H_{D'}=(1,h_1,h_2,1).
\]

Multiplication gives a pairing
\[
 (\n/\n^2)\times(\n^2/\n^3)\longrightarrow\n^3.
\]
It is nondegenerate on the second factor: a degree-two class annihilated by
\(\n\) would lift to a socle element and hence already lie in \(\n^3\).
Thus \(h_1\ge h_2\).  Choose the socle functional
\(\ell:\n^3\to\F_2\) with \(\ell(\theta)=1\).  The induced cubic functional
on \(\gr(D')\) pulls back along the surjection
\(\gr(B)\twoheadrightarrow\gr(D')\) to an element of \(U_B^\vee\) which
detects the universal target.  Since
\(\operatorname{Sym}^2(\n/\n^2)\twoheadrightarrow\n^2/\n^3\), its first
catalecticant factors as
\[
 \n/\n^2\twoheadrightarrow(\n^2/\n^3)^\vee
 \hookrightarrow\operatorname{Sym}^2(\n/\n^2)^\vee
\]
and has rank exactly \(h_2\).  The bounded MinRank UNSAT computation therefore
forces \(h_2\ge4\), conditional on its solver verdict.  Consequently
\[
 \length(D)\ge\length(D')=2+h_1+h_2\ge2+4+4=10.
\]
The ring \(R\) has Hilbert function \((1,4,4,1)\), so equality occurs.
\end{proof}

Because \(R\) is Gorenstein with socle \(\F_2\sigma\), every nonzero ideal
of \(R\) contains \(\sigma\).  Thus every proper quotient of this particular
base kills the displayed defect.  Moreover the Hopf coefficients contain
\(x,y,z,d\), which generate the entire ring, so the table is not secretly
written over a smaller coefficient subring.

The theorem concerns quotients of the cubic universal construction.  A
counterexample of length seven, eight, or nine could in principle arise at
greater Loewy depth, from a different closed fiber, or by a mechanism not
visible in this third-order chart.  Excluding that possibility requires a
separate global argument or search.

\section{Exact audit trail}\label{sec:audit}

Seven complementary checks were frozen in the workspace.
\begin{center}
\begin{longtable}{@{}p{0.49\textwidth}p{0.43\textwidth}@{}}
\toprule
file & role\\
\midrule
\endhead
\path{scripts/verify_rank4_length10_counterexample_20260710.py}
& dependency-free direct construction of \(R,A,\Delta,S,[4],[8]\), including
  exhaustive base socle and ring gates\\
\path{scripts/audit_rank4_length10_specialization_20260710.py}
& substitution of the \(45\) displayed values into the independently
  exported \(189\) equations and nine targets\\
\path{scripts/audit_rank4_psi_quasi_idempotent_20260710.py}
& exact filtered proof of \(\psi^2=2\psi\) on the universal cubic base\\
\path{scripts/audit_rank4_catalecticant_minimum_20260710.py}
& rank-four witness and bounded global rank-at-most-three UNSAT computation\\
\path{scripts/independent_verify_rank4_length10_counterexample_20260710.py}
& independent additive-coordinate model, exhaustive maximal-ideal powers,
  raw-equation evaluation, power maps, and convolution antipode\\
\path{scripts/audit_rank4_descriptions_20260710.py}
& exact regular \(\GL_4\) and reduced \(\GL_3\) matrices and inverses,
  recovery of all coordinate generators, dual Hopf table, and the basic
  rank-two subgroup\\
\path{scripts/audit_rank4_subgroups_deformations_20260711.py}
& exhaustive 1024-case classification of normalized rank-two Hopf quotients,
  adjoint nonnormality certificates, and exact length-\(6\to8\to10\)
  deformation tower\\
\bottomrule
\end{longtable}
\end{center}

The principal terminal output is:
\begin{verbatim}
BASE PASS: local characteristic 4; |R|=2^10; length 10
BASE PASS: Hilbert function (1,4,4,1); socle={0,s}; s=2Z!=0
BASE PASS: generator-level distributivity gates for the biadditive law
HOPF PASS: algebra, counit, multiplicativity, coassociativity
HOPF PASS: explicit two-sided convolution antipode constructed
POWER IDENTITY PASS: psi^2=2psi for psi=[2]^#-unit*epsilon
[4]^#(e1)=(s)e3
[4]^#(e2)=0
[4]^#(e3)=0
[8]^#=unit*epsilon PASS
NONCOMMUTATIVE PASS: Delta(e1) is not cocommutative
\end{verbatim}
The structural and deformation audit prints
\begin{verbatim}
GL4 PASS: displayed right-regular matrix and antipode inverse
GL3 PASS: quotient matrix invertible; e1,e2,e3 recovered exactly
DUAL HOPF PASS: all displayed products and coproducts
SUBGROUP CLASSIFICATION PASS: exactly 16 normalized Hopf quotients
ADJOINT PASS: H_x normal exactly after killing d
DEFORMATION PASS: lengths 6 -> 8 -> 10 with square-zero steps
DEFORMATION PASS: normality breaks at length 8; [4] at length 10
COMMUTATIVE QUOTIENT PASS: kill d and y+z (length 4)
\end{verbatim}
The raw universal export has SHA-256
\begin{verbatim}
686c35c91b74af2a414dd119024411266fb0056865dd1d19676c20485f93b6b7
\end{verbatim}
and the specialization audit prints
\begin{verbatim}
UNIVERSAL SPECIALIZATION PASS: 189/189 equations vanish
TARGETS PASS: [0,0,s,0,0,0,0,0,0], s=2Z!=0
\end{verbatim}

The universal quasi-idempotence check reproduces the filtered ranks
\[
 \operatorname{rank}_2=974,\qquad
 \operatorname{rank}_3=16787
\]
from \(45414\) cubic candidates, and reduces all nine coordinates of
\(T-2\phi\) to zero.  The MinRank script reconstructs all \(171\) cubic
relations before building its Boolean instance.

The direct verifier and raw-equation evaluator were also reconstructed
independently in a second implementation.  It exhaustively enumerated all
\(1024\) base elements, found socle \(\{0,\sigma\}\), evaluated all \(189\)
equations to zero, obtained the same target vector, and independently built
a two-sided convolution antipode.  Thus the headline result does not rest on
one parser, one Hopf implementation, or the MinRank solver.

\section{Conclusion}

The rank-four counterexample admits a human-sized form.  Its base is the
1024-element ring
\[
(\Z/4)[x,y,z,d]/
(2x,2y,2d,x^2-2,z^2-2,y^2,yz,zd,d^2-xd,yd-xz-xy),
\]
and its Hopf algebra is given by
\eqref{eq:A-multiplication}--\eqref{eq:Delta3}.  The fourth-power word is
nontrivial because
\[
 x^2(y+z)=2(y+z)=2z\ne0.
\]
This formula isolates the conceptual mechanism: the characteristic-two
binomial carry \(x^2=2\) multiplies a noncocommutative first-order direction,
and the result lands in the top socle.  In invariant language, the cubic
obstruction is \(\delta N\), and the exact endomorphism identity is
\(\psi^2=2\psi\).

The same group also has a concrete geometric and representation-theoretic
shape.  It is the relative complete intersection
\[
 R[U,V]/(U^2-xU-yV,V^2-zV),
\]
its dual measure algebra is the noncommutative PBW deformation in
\eqref{eq:dual-multiplication}, and its reduced regular representation is a
minimal closed immersion into \(\GL_{3,R}\).  It contains sixteen rank-two
Oort--Tate subgroups over \(R\), all isomorphic and all nonnormal.  Hence it
has subgroup chains but no normal composition filtration.

The deformation tower explains why.  Over \(R_6=R/(d)\), the group is a
genuine extension of two rank-two Oort--Tate groups.  The square-zero lift
to \(R_8\) turns on \(d\) and destroys normality, while preserving
\([4]=e\).  The final square-zero lift turns on
\(\sigma=dq=\omega d\), and the path
\(e_1\xmapsto{\omega}e_2\xmapsto{d}e_3\) produces
\([4]^\#(e_1)=\sigma e_3\).  Thus failure of a normal filtration and failure
of the order bound are successive, rather than simultaneous, deformation
obstructions.

The divided-power cubic
\(X^{[2]}Z+XYD+YD^{[2]}+Z^{[3]}\) packages the entire base multiplication.
Conditional on the recorded rank-at-most-three UNSAT verdict, its
catalecticant rank four is minimal in the full universal cubic dual space,
so length ten is the best possible compression by an Artin quotient of this
construction.  Determining whether an essentially
different counterexample exists over a base of length seven, eight, or nine
remains the next genuinely new minimization problem.

\appendix

\section{The complete universal specialization}\label{app:specialization}

Use the variable numbering of the universal chart:
\[
 p_{3\rho(i,j)+r-1}=m_{ij}^{\,r},\qquad
 p_{18+9(i-1)+3(j-1)+(k-1)}=c_{ijk},
\]
where \(\rho\) assigns the positions \(0,\ldots,5\), in order, to
\((1,1),(1,2),(1,3),(2,2),(2,3),(3,3)\).
The nonzero values are
\begin{center}
\begin{longtable}{@{}cl@{}}
\toprule
indices & value\\
\midrule
\endhead
\(p_0,p_8,p_{18},p_{38}\) & \(x\)\\
\(p_1\) & \(y\)\\
\(p_{10},p_{14}\) & \(z\)\\
\(p_{17}\) & \(\xi=xz\)\\
\(p_{21},p_{31}\) & \(y+z\)\\
\(p_{24},p_{32},p_{40},p_{44}\) & \(\xi+\omega=xz+xy=q\)\\
\(p_{30}\) & \(d\)\\
\(p_{39}\) & \(2+\eta=2+xd\)\\
\(p_{41}\) & \(\sigma=2z\)\\
\(p_{42}\) & \(x+d\)\\
\(p_{43}\) & \(y+z+\sigma\)\\
\bottomrule
\end{longtable}
\end{center}
All omitted \(p_i\) are zero.  Decoding this table gives exactly
\eqref{eq:A-multiplication}--\eqref{eq:Delta3}.

\section{The rank-four inverse-system map}

In the preliminary seven-variable compression, the surviving tangent
directions were labelled \(u,b,v,w,a,d,g\).  On these directions, the map
to \(\F_2\{X,Y,Z,D\}\) is
\[
 u,w\longmapsto X,\qquad b\longmapsto Y,
 \qquad v\longmapsto Z,
 \qquad a,g\longmapsto Y+Z,
 \qquad d\longmapsto D.
\]
In universal variable indices, its nonzero values are
\[
\begin{array}{c|c}
\{0,8,18,38\}&X\\
\{1\}&Y\\
\{10,14\}&Z\\
\{21,31,43\}&Y+Z\\
\{30\}&D\\
\{42\}&X+D.
\end{array}
\]
Pulling back the functional defined by \(\mathcal G\) annihilates the full
set of \(45414\) exhaustive cubic ideal candidates and evaluates to one on
the \(T_{13}\) normal class.  Its support in the \(15\)-variable cubic space
has size \(28\), and its catalecticant has rank four.

\end{document}
